% TOP_PROBLEM: {"problemId": "problem.banach-mazur-rotation-problem", "canonicalTitle": "Banach-Mazur Rotation Problem", "releaseRank": 249, "primaryDomain": "analysis_pde", "primaryDomainLabel": "Analysis & PDE", "displayStatus": "open", "releaseStatus": "open"}
% !TeX program = lualatex
% Definition and sources only; this file is not an LLM solution attempt.
\documentclass[11pt]{article}
\usepackage[margin=25mm]{geometry}
\usepackage{fontspec}
\setmainfont{Latin Modern Roman}
\usepackage{amsmath,amssymb,mathtools}
\usepackage{unicode-math}
\setmathfont{Latin Modern Math}
\usepackage{xurl,hyperref}
\hypersetup{colorlinks=true,urlcolor=blue}
\setcounter{secnumdepth}{0}
\setlength{\parindent}{0pt}
\setlength{\parskip}{5pt}
\setlength{\emergencystretch}{3em}
\begin{document}
\section*{\texorpdfstring{Banach-Mazur Rotation Problem}{Banach-Mazur Rotation Problem}}\label{249}
\subsection{Definitions and mathematical statement}
A real Banach space is a complete normed real vector space. Its surjective linear isometries are linear bijections $U:X\to X$ with $\|Ux\|=\|x\|$. Transitivity on the unit sphere means
\[
 \forall x,y\in X,\quad\|x\|=\|y\|=1
 \quad\Longrightarrow\quad\exists U\text{ as above with }Ux=y.
\]
For infinite-dimensional separable $X$, does this force the norm to arise from an inner product, equivalently the parallelogram identity
$\|x+y\|^2+\|x-y\|^2=2\|x\|^2+2\|y\|^2$?
The conclusion is linear \emph{isometry} to a Hilbert space, not merely a bounded linear isomorphism. Approximate transitivity is not the selected hypothesis.
\par\smallskip\textit{Formulation and concept references:} \href{https://arxiv.org/abs/2012.08344}{[S1]}.

\subsection{Short English statement}
If every unit vector in a separable real Banach space can be carried exactly to every other by a linear isometry, must its norm be Hilbertian?
\subsection{Sources}
\begin{itemize}
\item[S1]On Mazur rotations problem and its multidimensional versions.
\url{https://arxiv.org/abs/2012.08344}.
\textit{Formulation source.}
\end{itemize}
\end{document}
